\chapter{总结与展望}
本文立足于科学计算与机器学习的交叉领域，以复杂的偏微分系统为研究对象，系统探讨了基于深度神经网络的建模与计算方法。针对传统网格类数值方法在高维空间中面临的”维数灾难”，以及现有深度学习求解器在处理复杂物理约束和病态问题时的理论空白，本文秉持算法设计适配复杂物理场景、理论分析提供定量保障、数值实验验证实际泛化能力的研究主线，取得了一系列具有理论价值和应用前景的研究成果。全文的主要工作总结如下：
\begin{itemize}
    \item 在高维椭圆型多重特征值问题方面（第三章）：针对现有无约束损失函数易陷入模式崩溃的难题，本文提出了一种基于惩罚变分原理的 DRM 求解框架。通过引入软边界惩罚与正交化约束，实现了多重激发态的同步求解。在理论上，本文建立了误差分解体系（逼近误差、累积误差与统计误差），并推导了由网络结构与样本量显式表达的误差上界，在数学上表明该算法具有缓解维数灾难的潜力。
    \item 在高维半线性椭圆偏微分方程方面（第四章）：面对强非线性特征导致的优化困难，本文在次临界与临界增长假设下，构造了基于修正能量泛函的极小化算法。本文创新性地建立了一套包含截断界的误差分解框架，推导了 $H^1$ 范数下的总误差上界，并给出了算法在求解 Robin 和 Dirichlet 问题时的渐近收敛率。针对多种类型非线性 PDE 问题的数值实验，验证了该算法处理高维复杂非线性特征时的精度与鲁棒性。
    \item 在椭圆偏微分方程反势能问题方面（第五章）：针对反问题固有的病态性及”半收敛”现象，本文提出了一种融合 PINNs 与 Tikhonov 正则化（$H^2$ 范数惩罚）的重构框架。本文建立了该神经网络格式的条件稳定性估计，并推导了重构误差关于观测噪声 $\delta$ 的渐近收敛率。肿瘤热疗及污染物扩散等反演实验表明，该方法在保持较高精度的同时，对观测噪声展现出较好的鲁棒性。
\end{itemize}

尽管本文在基于神经网络求解复杂偏微分系统方面取得了一定的理论与算法进展，但受限于深度学习优化理论的复杂性以及 PDE 本身的物理特性，本领域仍有诸多具有挑战性的方向有待探索。基于本文的研究基础，未来的工作可从以下几个维度展开：
\begin{itemize}
    \item 高阶特征值算法的拓展与优化：本文提出的基于正交惩罚的变分算法在求解前几阶特征值时表现良好。然而，当目标特征值的阶数 $k$ 进一步增大时，正交约束项的数量将线性增加，这不仅会导致损失函数的优化变得困难，还可能引发累积误差放大效应。未来的研究可以考虑引入不完全正交化策略、子空间迭代思想或基于复流形优化的神经网络结构，以进一步提升算法对高阶特征值（乃至连续谱带）的适用性与计算稳定性。
    \item 弱化非线性偏微分方程的理论假设：本文在第三章对半线性 PDE 的理论分析建立在 Sobolev 次临界/临界增长以及势函数单调性等假设之上，这主要为了保证能量泛函的严格凸性与弱解的唯一性。在未来的工作中，可以将算法与理论拓展至更为宽松的非线性场景中。
    \item 探索多源正则化机制在反演网络中的应用：在第五章的 Tikhonov-PINNs 框架中，本文主要采用了 $H^2$ 范数进行正则化，这对于平滑势函数表现出较好的理论性质。未来可以探讨在 PINNs 损失中引入全变差正则化、稀疏 $L^1$ 正则化等正则化项。如何为这些非光滑正则化项建立误差分析与稳定性界限，将是一个具有理论价值的课题。
    \item 面向稀疏/局部观测数据的反演算法研究：本文当前的反问题理论分析与数值实验均建立在对全区域 $\Omega$ 内状态变量 $u$ 进行全场观测的假设之上。然而，在真实的工程应用中，通常只能获取有限的局部内部传感器数据或纯边界测量数据。未来需要将本文的算法框架推广至局部观测场景。在理论层面，这要求结合偏微分方程的性质建立新的条件稳定性界；在算法层面，则需要设计更高效的采样策略或注意力机制，以克服因观测数据稀疏带来的不适定性。
\end{itemize}